\section{Conclusion}
\label{sec:conclusion}

We have presented a cybernetic model of socialist economic planning
that integrates dynamic input-output analysis, endogenous growth
inference, and shadow price feedback.
The key contributions are:
(i)~embedding investment into the technical structure via the
growth-augmented matrix $\bar{A} = A + \delta B$;
(ii)~inferring growth endogenously from revealed demand rather than
imposing it exogenously;
(iii)~deriving shadow prices from constrained welfare maximization to
guide decentralized coordination;
and (iv)~implementing adaptive preference learning to enable
demand-led structural change.

This framework takes steps toward addressing longstanding critiques of central planning \citep{hayek1945use, kornai1992socialist} by incorporating price signals, distributed information, and adaptive feedback into the planning process. Our benchmarking against a Heterogeneous Agent New Keynesian (HANK) baseline demonstrates that the cybernetic system achieves an 8.3\% median welfare advantage by optimizing allocative efficiency, even while producing lower raw growth volume (12.3\% vs 15.5\%) than the uncoordinated market. The model similarly abstracts from the instability and distributional dynamics documented in market economies \citep{minsky1986stabilizing, kalecki1971selected}, though the statistical evidence indicates that coordinated planning can significantly reduce the utility losses associated with monetary distortions and uninsurable liquidity risk. Modern computational capabilities make cybernetic planning systems of this kind practically implementable in ways unavailable to earlier theorists, building on the foundations laid by \citet{lange1936economic}, \citet{kantorovich1965best}, and \citet{leontief1986input}, though the gap between a calibrated simulation on national accounts data and a functioning planning system for a real economy should not be understated.